\section{Nine-Point Enclosure Obstruction}
\label{sec:strategy4-reader}

Assume throughout that all six vertex triangles are $\Vdzero$, that they
cover $\partial H$, and that exactly one is supercritical.  Rotate its index
to $4$.  Lemma~\ref{lem:ab-extreme-jump} gives strict handoff parameters
\begin{equation}
 \xi_3<\xi_2<\xi_1<\xi_0<\xi_5<\xi_4,
 \label{eq:reader-extreme-chain}
\end{equation}
and
\begin{equation}
 a=1-\xi_3,
 \qquad
 b=\xi_4,
 \qquad
 p=1-b,
 \qquad
 q=1-a,
 \label{eq:reader-ab-parameters}
\end{equation}
with
\begin{equation}
 0<a,b<1,
 \qquad
 a+b>1,
 \qquad
 a^2+ab+b^2<1.
 \label{eq:reader-ab-domain}
\end{equation}
Every selected backward reach is at least $p$, and every selected forward
reach is at least $q$.

\subsection{Six radial points and three frontier points}

Let
\[
 P_A=V_4+a(V_3-V_4),
 \qquad
 P_B=V_4+b(V_5-V_4),
\]
and define the feasible $AB$-union
\[
 \mathcal R_{AB}(a,b)=
 \bigcup\left\{T\cap H:
 \begin{array}{l}
 T\text{ is a closed unit equilateral triangle},\\
 V_4,P_A,P_B\in T
 \end{array}\right\}.
\]
For $(x,y)\in[0,1]^2$, put
\[
 c_{\max}(x,y)=\max\{c:(x,y,c)\in\mathcal A\},
\]
and define
\[
 c_\ast=c_{\max}(p,q),
 \qquad
 h=\frac{\sqrt3}{2},
 \qquad
 r_\ast=h(1-c_\ast).
\]
The six radial points
\[
 P_i^{\rm rad}=(1-c_\ast)V_i,
 \qquad 0\le i\le5,
\]
are excluded from every vertex triangle.  Therefore they lie in $U_C$, and
convexity gives
\begin{equation}
 \mathbb B(0,r_\ast):=\{X:\lVert X\rVert\le r_\ast\}\subset U_C.
 \label{eq:reader-forced-disk}
\end{equation}

The boundary of $\mathcal R_{AB}(a,b)$ at the supercritical corner has two
circle pieces and two line segments.  Denote the two line segments by
$\ell_-$ and $\ell_+$ in their cyclic order, and let $\mathcal C_-$ and
$\mathcal C_+$ be the two adjacent-handoff unit circles specified in
Lemma~\ref{lem:fixed-line-signs}.  Define
\[
 Q_0=\ell_-\cap\ell_+,
\]
and let $Q_-$ and $Q_+$ be the unique intersections
\[
 Q_-\in\ell_-\cap\mathcal C_-,
 \qquad
 Q_+\in\ell_+\cap\mathcal C_+
\]
on the corresponding frontier segments.  The fixed-line sign calculation
proves uniqueness and excludes all three points from the two adjacent vertex
triangles.  Distance inequalities exclude the other three nonsupercritical
vertex triangles, and frontier membership excludes $T_4$.  Hence
\begin{equation}
 Q_-,Q_0,Q_+\in U_C.
 \label{eq:reader-forced-asymmetric}
\end{equation}
The nine forced points are the six $P_i^{\rm rad}$ and the three points
$Q_-,Q_0,Q_+$.  For the enclosure argument we use the compact set
\begin{equation}
 K_{\rm wit}(a,b)=\mathbb B(0,r_\ast)\cup\{Q_-,Q_0,Q_+\}
 \subset U_C.
 \label{eq:reader-witness-set}
\end{equation}

\subsection{Rational inner approximants}

The coordinates of $Q_-$ and $Q_+$ contain two independent quadratic
radicals.  Lemma~\ref{lem:technical-newton-reduction} constructs points
\[
 P_-\in(Q_0,Q_-),
 \qquad
 P_0=Q_0,
 \qquad
 P_+\in(Q_0,Q_+)
\]
whose coordinates are rational in $a,b$ and the single radical
\[
 \Delta=\sqrt{4(a^2+ab+b^2)-3}.
\]
These rational inner approximants satisfy
\begin{equation}
 \widehat K:=\conv\bigl(\mathbb B(0,r_\ast)\cup\{P_-,P_0,P_+\}\bigr)
 \subseteq\conv(K_{\rm wit}),
 \qquad
 \Lambda(K_{\rm wit})\ge\Lambda(\widehat K).
 \label{eq:reader-newton-reduction}
\end{equation}
The appendix obtains $P_-$ and $P_+$ by one inward Newton step on the two
frontier lines.

\subsection{A cyclic support-arc covering lemma}

For $X\in\mathbb R^2$ and $r\ge0$, write
\[
 \mathbb B(X,r)=\{Y:\lVert Y-X\rVert\le r\},
\]
and define the level-$h$ support arc
\[
 \operatorname{Cap}(X,r)
 =\{n\in\mathbb S^1:\langle X,n\rangle+r\ge h\}.
\]

\begin{lemma}[Cyclic support-arc covering]
\label{lem:reader-cap-chain}
Let $K$ be compact and put
\[
 \Sigma_K=K+\mathsf RK+\mathsf R^2K.
\]
Suppose $\Sigma_K$ contains the full $\mathsf R$-orbits of
\[
 \mathbb B(P_-,2r_\ast),
 \quad
 \mathbb B(P_0,2r_\ast),
 \quad
 \mathbb B(P_+,2r_\ast),
 \quad
 \mathbb B(P_\ast,r_\ast),
 \qquad
 P_\ast=P_++\mathsf RP_-.
\]
Let $\theta_X$ denote the argument of the ray through $X$.  Assume the rays
\[
 P_-,P_0,P_+,P_\ast,\mathsf RP_-
\]
occur in this cyclic order, every corresponding support arc is connected with
half-width less than $\pi/2$, and for each consecutive pair $(X,r),(Y,s)$ in
\[
 (P_-,2r_\ast),
 (P_0,2r_\ast),
 (P_+,2r_\ast),
 (P_\ast,r_\ast),
 (\mathsf RP_-,2r_\ast)
\]
one has
\[
 \operatorname{Cap}(X,r)\cap\operatorname{Cap}(Y,s)
 \cap\{n:\arg n\in[\theta_X,\theta_Y]\}\ne\varnothing,
\]
where $[\theta_X,\theta_Y]$ is the intervening short angular interval.  Then
\[
 \Lambda(K)\ge1.
\]
\end{lemma}

\begin{proof}
The four consecutive intersections and cyclic order make the five support
arcs cover the sector from the ray of $P_-$ to the ray of $\mathsf RP_-$.  Their
$\mathsf R$-orbits cover $\mathbb S^1$.  Hence, for every unit vector $n$, one
of the displayed disks has support at least $h$ in direction $n$.  Thus
\[
 h_{\Sigma_K}(n)\ge h.
\]
Since
\[
 h_{\Sigma_K}(n)=\sum_{j=0}^2h_K(\mathsf R^jn),
\]
Proposition~\ref{prop:universal-enclosure-gauge} gives $\Lambda(K)\ge1$.
\end{proof}

\begin{proposition}[Four support-arc intersections]
\label{prop:reader-four-overlaps}
Assume \eqref{eq:reader-ab-domain} and $c_\ast>2/3$.  The rational inner
approximants satisfy
\[
 \lVert P_-\rVert,\lVert P_+\rVert>r_\ast,
\]
the five rays in Lemma~\ref{lem:reader-cap-chain} occur in the required cyclic
order, and all four consecutive support-arc intersections in that lemma are
nonempty.
\end{proposition}

\begin{proof}
Proposition~\ref{prop:technical-four-overlaps} proves the ray order, radius
bounds, and the two equal-radius intersections.  Theorem
\ref{thm:paper-exact-mixed-certificate} proves the two unequal-radius
intersections by exact Gram reductions and Bernstein-positivity certificates.
\end{proof}

\subsection{Enclosure contradiction}

\begin{theorem}[Witness enclosure]
\label{thm:reader-witness-enclosure}
For every pair $(a,b)$ satisfying \eqref{eq:reader-ab-domain},
\[
 \Lambda(K_{\rm wit}(a,b))\ge1.
\]
\end{theorem}

\begin{proof}
The centered disk alone gives
\[
 \Lambda(K_{\rm wit})\ge\frac{3r_\ast}{h}=3(1-c_\ast),
\]
so the result follows immediately when $c_\ast\le2/3$.  For $c_\ast>2/3$,
apply \eqref{eq:reader-newton-reduction},
Proposition~\ref{prop:reader-four-overlaps}, and
Lemma~\ref{lem:reader-cap-chain}.
\end{proof}

\begin{theorem}[Zero-gap obstruction]
\label{thm:reader-zero-gap-obstruction}
There is no cover of $H$ by open unit equilateral triangles
$U_C,U_0,\ldots,U_5$ such that every $U_i$ is $\Vdzero$, the six vertex
triangles cover $\partial H$, and exactly one vertex triangle is
supercritical.
\end{theorem}

\begin{proof}
The nine-point argument gives the compact containment
\[
 K_{\rm wit}\subset U_C,
\]
whereas Theorem~\ref{thm:reader-witness-enclosure} gives
$\Lambda(K_{\rm wit})\ge1$.  Every compact subset of an open unit equilateral
triangle is contained in a closed equilateral triangle of side strictly below
one, which would imply $\Lambda(K_{\rm wit})<1$.
\end{proof}

\begin{proposition}[Cases excluded by the nine-point obstruction]
\label{prop:reader-ab-core-branches}
Theorem~\ref{thm:reader-zero-gap-obstruction} excludes:
\begin{enumerate}
\item the $\CEzero$, $N_+=1$, all-$\Vdzero$ case;
\item the no-gap $\CEone$ and $\CEtwo$, $N_+=1$, all-$\Vdzero$ cases.
\end{enumerate}
\end{proposition}

\begin{proof}
In the CE0 case, the six vertex triangles must cover the perimeter, since any
missed boundary point would lie in the open center triangle and create a
positive center trace.  In the CE1 and CE2 cases, absence of a boundary gap
means that the two incident vertex traces cover every edge.  The
unique-supercritical hypothesis is common to all three cases.
\end{proof}
